\section{Tìm kiếm A* (A-star Search)}
\label{sec:bloxorz-astar}

A* sử dụng cùng chi phí tại Công thức~\eqref{eq:bloxorz-action-cost}, ưu tiên
$f(s)=g(s)+h(s)$, với $h$ là hàm ước lượng (heuristic). Lớp \texttt{AStarFrontier} kế thừa frontier chi phí của UCS và chỉ đổi
hàm ưu tiên; cơ chế cải thiện đường đi và mở lại trạng thái được dùng chung, bảo đảm một trạng thái từng gặp vẫn có thể
nhận đường rẻ hơn. Các bộ giải giữ nguyên giao diện
\texttt{solve(stage, caps)} và kiểu kết quả đã đóng băng.

\subsection{Thiết kế hàm ước lượng}
Trên màn không có công tắc tách, đặt:
\begin{align}
  d(s)&=\min_{(r,c)\in O(s)}\bigl(|r-r_G|+|c-c_G|\bigr),\\
  h(s)&=\frac{c_{\min}\,d(s)}{2}.
  \label{eq:bloxorz-heuristic}
\end{align}
Khoảng cách được lấy từ ô bị chiếm gần đích nhất, không từ một tọa độ tham
chiếu có thể làm quá cận. Dữ liệu tĩnh được kiểm tra một lần để tạo hàm ước
lượng, không quét lưới trên từng nút.

\subsection{Không đánh giá quá cao và tính nhất quán}
\textbf{Không đánh giá quá cao (admissibility).}
Một lần lăn giảm $d$ nhiều nhất hai ô. Bỏ qua hướng khối, vật cản, cầu và ô dễ
vỡ chỉ làm bài toán dễ hơn, nên cần ít nhất $d(s)/2$ bước để tới đích. Mỗi bước
tốn ít nhất $c_{\min}$, suy ra $h(s)\leq h^*(s)$, với $h^*$ là chi phí tối ưu
còn lại. Lập luận dùng cận chi phí của chính mô hình không đồng nhất đang chạy.

\textbf{Nhất quán (consistency).}
Với mỗi cạnh không dịch chuyển, $d(s)\leq d(s')+2$, do đó:
\begin{equation}
  h(s)\leq h(s')+c_{\min}\leq h(s')+c(s,a,s').
\end{equation}
Tại đích, $h=0$. Mọi phụ phí đều không âm nên bất đẳng thức vẫn đúng trên
các bước có tiếp xúc ô dễ vỡ.

\subsection{Giới hạn do dịch chuyển và hướng khối}
Công tắc tách có thể dịch chuyển xa tùy ý, phá vỡ giới hạn hai ô. Vì vậy dùng
$h=0$ trên \emph{mọi trạng thái của màn có công tắc tách}, kể cả trước khi tách.
Phương án này không đánh giá quá cao và nhất quán, nhưng làm A* trở thành UCS.
Chỉ tắt hàm ước lượng sau khi khối đã tách là không đủ vì chính cạnh kích hoạt
dịch chuyển cũng có thể làm giảm khoảng cách rất lớn.

Không cộng tùy ý một đơn vị phạt khi khối nằm. Nếu khối nằm ngay cạnh đích và
có thể đứng lên đích trong một bước chi phí $1$, khoảng cách ô gần nhất là $1$,
nên $h=1/2$. Cộng thêm $1$ cho hướng nằm tạo ước lượng $3/2>1$, vượt chi phí
thật. Vì vậy phần cài đặt không sử dụng hạng phạt hướng chưa được chứng minh.

\subsection{Kiểm chứng và phạm vi kết luận}
Kiểm thử phát lại đường đi, so sánh chi phí UCS/A*, kiểm tra giới hạn, trường
hợp không có nghiệm, bản ghi cũ và mở lại nút đã mở rộng dưới một hàm ước lượng
không đánh giá quá cao nhưng không nhất quán. Trên ba màn nhỏ có đáp án tính
tay, chi phí còn lại được tính độc lập ở mọi trạng thái có đường tới đích và
đối chiếu với $h$. Kiểm tra mở rộng trên mười màn chính thức duyệt $2\,946$
trạng thái đạt được, tạo đồ thị cạnh ngược để tính chi phí tối ưu còn lại và
kiểm tra tính nhất quán trên mọi cạnh.

Ví dụ kiểm tra tính đúng trên L2-02: phép duyệt ít bước nhất chỉ nằm trong
kiểm thử tìm đường $46$ bước, chi phí $64$; UCS/A* chọn đường $49$ bước, chi
phí $58$. Đây là bằng chứng cho sự khác nhau giữa mục tiêu số bước và mục tiêu
chi phí, không phải số liệu đo thời gian hoặc bộ nhớ. Các chứng minh ở trên
áp dụng cho mô hình đã nêu; phép duyệt hữu hạn là kiểm chứng bổ sung, không
thay thế chứng minh cho mọi màn có thể được tạo.

Với hàm ước lượng nhất quán đang dùng, mỗi trạng thái được mở rộng nhiều nhất
một lần hợp lệ. Gọi $V,E$ là số trạng thái/cạnh đạt được, thời gian dùng hàng
đợi ưu tiên là $O((V+E)\log(V+E))$, bộ nhớ $O(V+E)$ kể cả bản ghi thay thế.
Không suy ra A* luôn nhanh hơn UCS: độ chặt của cận và chi phí tính hàm ước
lượng đều ảnh hưởng kết quả. Trên màn tách, hai thuật toán có cùng chiến lược
vì $h=0$. Đánh giá hiệu năng toàn nhóm chờ bộ công cụ đo đạc và phiên bản tích
hợp cuối; phần này không đưa số đo chưa được thực hiện.
